\section{Results}
\label{sec:results}

% EIXO (b), decisao de 11/09: o achado de informacao (H11) abre a secao; os
% efeitos sobre tempos, comercio e eficiencia vem como contexto — tres
% resultados nulos honestos e um efeito documental estreito.

% ATENCAO (23/09): o salto de cobertura e essencialmente Santos (0 -> 67% na janela);
% o efeito medio nao e robusto (DiD A 0,197 p_wcb 0,38; com tendencia por porto ~0).
% Ver tab:pretrend e outputs/tables/coverage_jump_by_port.csv (decisao 42).
\subsection{What the single window changed first: the information the regulator receives}
\label{sec:res_info}

Before the single window, the post-operation waiting time ($T_4$) was
recorded for \CobSantosPre\ of vessel calls in Santos and \CobRioPre\ in Rio
de Janeiro; after it, for \CobSantosPos\ and \CobRioPos\ respectively
(\Cref{fig:coverage}). \Cref{tab:covefeito} reports the corresponding
difference-in-differences estimates.

\input{tables/tab06_cobertura_efeito}

\begin{figure}[htbp]
\centering
\includegraphics[width=\textwidth]{fig13_es_coverage.pdf}
\caption{Event study of the probability that $T_4$ is recorded, dated public ports against
never-treated private terminals, 2010--2015. Sun--Abraham estimator with port and
month fixed effects; month $-1$ is the reference. Shaded band: 95\% uniform
(sup-$t$) band \citep{montielolea2019}; vertical lines: pointwise 95\% intervals; standard errors
clustered by port. Source: ANTAQ waterway statistics, own elaboration.}
\label{fig:escoverage}
\end{figure}

\PROSA

\subsection{Operational times: no aggregate effect}
\label{sec:res_times}

Across dated public ports, the static effect of the single window on
documentary waiting times is \TdoisBetaFull\ for $T_2$ (wild-bootstrap
$p=\TdoisPwcbFull$) and \TquatroBetaFull\ for $T_4$ ($p=\TquatroPwcbFull$);
the physical queue ($T_1$) shows no movement (\Cref{tab:inferencia}).

\input{tables/tab07_inferencia}

\begin{figure}[htbp]
\centering
\includegraphics[width=\textwidth]{fig02_coortes_psp.pdf}
\caption{Documentary waiting times by adoption cohort, quarterly medians,
2010--2015. Dashed lines mark the first production date of each cohort.
Raw medians mix changes in times with changes in which calls are recorded
(\Cref{sec:res_info}); descriptive evidence only. Source: ANTAQ waterway statistics, own elaboration.}
\label{fig:cohorts}
\end{figure}

\PROSA

% ATENCAO (23/09, decisao do autor pendente): o titulo e o paragrafo abaixo
% foram escritos antes do teste de pre-tendencias. Com tendencia linear por porto
% e com HonestDiD (M=0) o efeito da cabotagem NAO sobrevive (tab:pretrend,
% fig:honestdid, docs/decisions_log.md decisao 42). Reescrever ao redigir a prosa.
\subsection{Where a documentary effect survives: recurrent cabotage calls}
\label{sec:res_nav}

The only estimate that survives both few-cluster inference and the
anticipation placebo is the post-operation waiting time of recurrent
cabotage vessels: \CabRecTquatroBeta\ in logs (about \CabRecTquatroPct\
shorter, wild-bootstrap $p=\CabRecTquatroPwcb$; \Cref{tab:heterog}).

\input{tables/tab08_heterogeneidade}

\begin{figure}[htbp]
\centering
\includegraphics[width=\textwidth]{fig12_es_cabotage.pdf}
\caption{Event study of documentary waiting times for recurrent cabotage vessels (three or
more calls at the same port in the previous twelve months), design A
(not-yet-treated public ports, 2010--2013). Sun--Abraham estimator with port,
month and navigation fixed effects and cargo controls; month $-1$ is the
reference. Shaded band: 95\% uniform (sup-$t$) band \citep{montielolea2019}; vertical lines: pointwise
95\% intervals; standard errors clustered by port. Source: ANTAQ waterway statistics, own elaboration.}
\label{fig:escabotage}
\end{figure}

\PROSA

\subsection{Long-run dynamics, 2010--2026}
\label{sec:res_long}

\begin{figure}[htbp]
\centering
\includegraphics[width=\textwidth]{fig08_horizonte_longo.pdf}
\caption{Long-horizon event study for documentary waiting times, dated public
ports against never-treated private terminals, 2010--2026. Sun--Abraham
estimator with port, month and navigation fixed effects and cargo controls.
Shaded band: 95\% uniform (sup-$t$) band \citep{montielolea2019}; vertical lines: pointwise 95\%
intervals; standard errors clustered by port. Source: ANTAQ waterway statistics, own elaboration.}
\label{fig:longrun}
\end{figure}

\PROSA

\FloatBarrier
